%! Project: Design and Conduct a Survey — Unit 1, Lesson 1.8 — Warm-Up (Answer Key)
% ONE PAGE, blank and key. Three items at 12pt (the stats-model size), \workrowsep 50pt,
% item spacing 22pt, so each item gets real handwriting room and still fits. Do not add a
% fourth item — the page is full.
% GRADUAL RELEASE: the warm-up activates prior knowledge before the guided notes. Items 1
% and 2 are the arithmetic spine of the whole project, run once on a study that is already
% finished — the Student Council's stratified 60 — so the numbers the notes open with
% (35%, 315) are already on the page. Item 3 separates the two groups WITHOUT the
% vocabulary and lands the 100% check that row 2 names as instrument rule 3.
\documentclass[12pt]{article}
\usepackage{saar-article}
\usepackage{saar-key}   % pulls in -boxes; do NOT also load -boxes
\newcommand{\TallMath}[1]{$\displaystyle #1\rule[-1.4em]{0pt}{3.2em}$}
% Word bank strip for every fill-in cluster (user request 2026-09-06). Defined per document;
% the key inherits it and prints the same strip, so heights match.
\newcommand{\wordbank}[1]{\par\vspace{2pt}\noindent\fcolorbox{goldacc}{goldbg}{\parbox{\dimexpr\linewidth-2\fboxsep-2\fboxrule\relax}{\small\textbf{Word bank:} #1}}\par\vspace{4pt}}
% Handwriting room inside every \begin{work} block. Moves the blank and the
% key together, so the two can never drift apart.
\setlength{\workrowsep}{50pt}

\begin{document}

\pageheader{Unit 1, Lesson 1.8}{Warm-Up --- Answer Key}
% NAMESTRIP: no \namedateperiod here — mirrors the blank. NO teachernote here —
% teacher prose lives in the lesson plan.

\vspace{0.15in}

\begin{notesbox}{Spiral Review --- A Study That Is Already Finished}
\normalsize
The Riverbend Student Council has a \$$3{,}000$ grant and asked the school which
one thing it should fund. It drew a \textbf{stratified sample of $60$} of the
$900$ students and got one answer from each: club funding $21$, field trips
$18$, gym equipment $15$, longer lunch $6$.

\begin{enumerate}[leftmargin=1.8em, itemsep=22pt, topsep=10pt]

  \item What \textbf{percent} of the $60$ students asked chose club funding?

        \begin{work}
          \dfrac{21}{60} &= 0.35 \\
                         &= 35\%
        \end{work}

  \item Based on that survey, \textbf{predict} how many of all $900$ Riverbend
        students would choose club funding.

        \begin{work}
          0.35 \times 900 &= 315 \text{ students}
        \end{work}

  \item Two different groups of students appear in the paragraph above. Write how many are in each group.
        \wordbank{60 \;\textbullet\; 900 \;\textbullet\; 315 \;\textbullet\; exactly one \;\textbullet\; more than one}

        \vspace{4pt}
        \renewcommand{\arraystretch}{2.6}
        \begin{tabularx}{\linewidth}{@{} X p{3.4cm} @{}}
        \toprule
        \textbf{Group} & \textbf{How many} \\
        \midrule
        The group the council actually collected data from & \ans{60} \\
        The whole group its report is about                 & \ans{900} \\
        \bottomrule
        \end{tabularx}

        \vspace{10pt}
        The four percents ($35$, $30$, $25$, $10$) add to $100\%$. How many
        choices did each student pick? \ans{exactly one}

\end{enumerate}
\end{notesbox}

\end{document}
